% Generated draft: requires mathematical and source review.
\documentclass[11pt]{article}
\usepackage[margin=25mm]{geometry}
\usepackage{fontspec}
\setmainfont{Latin Modern Roman}
\usepackage{amsmath,amssymb,mathtools,unicode-math}
\setmathfont{Latin Modern Math}
\usepackage{enumitem,needspace,xurl,hyperref}
\setlength{\parindent}{0pt}
\setlength{\parskip}{5pt}
\setlength{\emergencystretch}{3em}
\newcommand{\N}{\mathbb{N}}
\newcommand{\Z}{\mathbb{Z}}
\newcommand{\Q}{\mathbb{Q}}
\newcommand{\R}{\mathbb{R}}
\newcommand{\C}{\mathbb{C}}
\newcommand{\F}{\mathbb{F}}
\newcommand{\A}{\mathbb{A}}
\newcommand{\PP}{\mathbb{P}}
\newcommand{\E}{\mathbb{E}}
\newcommand{\eps}{\varepsilon}
\newcommand{\dd}{\,\mathrm{d}}
\begin{document}
\section*{UnsolvedMath ID: 30000021}
\textbf{Source problem code:} catalogue entry
\subsection{Definitions and mathematical statement}
The intended sphere is apparently the two-dimensional unit sphere, but its dimension is not stated in the supplied question. Under this reading let
\[
S^2=\{x\in\R^3:\|x\|_2=1\},\qquad \sigma(S^2)=4\pi,
\]
where \(\sigma\) is surface area measure. A spherical \(n\)-design is a finite set \(X\subset S^2\) of distinct points such that
\[
\frac{1}{|X|}\sum_{x\in X}p(x)=\frac{1}{4\pi}\int_{S^2}p(x)\,\dd\sigma(x)
\quad\text{for every }p\in\R[x_1,x_2,x_3]\text{ of total degree at most }n.
\]
The source question, on this interpretation, is whether
\[
\forall n\in\Z_{\geq0}\quad\exists X_n\subset S^2\quad
|X_n|=(n+1)^2\quad\text{and}\quad X_n\text{ is a spherical }n\text{-design}.
\]
Thus equal weights and exactly \((n+1)^2\) distinct nodes are required; neither arbitrary positive cubature weights nor a bound with an unspecified constant addresses this precise target. ``Degree \(n\)'' is read in the usual design sense of all degrees at most \(n\). The case \(n=0\) is harmless if the intended indexing begins at \(n=1\).

The supplied literature triage dated 2026-08-21 reports rigorous computational existence verification through degree 100 and numerical evidence to substantially higher degrees, while reporting that no proof for every degree was located. These finite-degree and numerical results do not settle the universal quantifier in the displayed question; their exact ranges and hypotheses require independent checking in the supplied references.
\subsection{Short English statement}
On the intended \(S^2\) interpretation, does every degree \(n\) admit an equal-weight spherical design with exactly \((n+1)^2\) distinct points? The supplied review reports finite-degree computational verification, not a proof for all degrees.
\subsection{Sources}
\begin{itemize}
\item Local UnsolvedMath registry, permanent ID 30000021, catalogue entry, \emph{Spherical Designs with Square Cardinality}.
\item \emph{Geomathematik} (2005), Oberwolfach Reports, pp. 1371--1448, DOI \url{https://doi.org/10.4171/owr/2004/27}.
\item Chen--Frommer--Lang, supplied DOI \url{https://doi.org/10.1007/s00211-010-0332-5}; no article title or theorem number is included in the record.
\item Supplied reference labeled ``spherical design verification'': \url{https://arxiv.org/abs/1502.03562}.
\item Supplied dated literature assessment: checked 2026-08-21; open status and verification ranges are reported from that curation.
\end{itemize}
\end{document}
